\chapter{Gestion de version}\label{chap:git}

\section{Principes}
Le dépôt \code{jee-ecommerce-jpa-mvc2} est hébergé sur GitHub. Trois règles ont été suivies.

\begin{itemize}
  \item \textbf{Commits atomiques} : un commit correspond à un changement cohérent (une couche, une fonctionnalité), qui
        compile seul. L'historique se lit alors comme le déroulement du projet.
  \item \textbf{Messages normalisés} : la convention \textit{Conventional Commits} \cite{conventional}, avec les types
        \code{feat}, \code{fix}, \code{test}, \code{docs}, \code{chore}, et une portée (\code{persistence}, \code{service},
        \code{web}, \code{views}).
  \item \textbf{Rien de secret dans le dépôt} : \code{db.properties} est dans \code{.gitignore}, seul le modèle
        \code{db.properties.example} est versionné. Avant chaque commit, \code{git status} permet de vérifier que le fichier n'apparaît pas.
\end{itemize}

\section{Historique}
Le tableau~\ref{tab:commits} donne la séquence des commits, qui suit l'ordre des étapes.

\begin{table}[htbp]
\centering
\caption{Séquence de commits}\label{tab:commits}
\small
\begin{tabularx}{\linewidth}{@{}lX@{}}
\toprule
\textbf{Étape} & \textbf{Message}\\
\midrule
Cadrage & \code{docs: add class diagram and project scope}\\
Maven & \code{chore: add Maven project and gitignore}\\
Persistance & \code{feat(persistence): add JPA entities, persistence unit and schema tool}\\
Schéma & \code{fix(persistence): require a role on every internaute}\\
Services & \code{feat(service): add business services, CDI wiring and transaction runner}\\
Contrôleurs & \code{feat(web): add MVC2 servlets, filters and security}\\
Vues & \code{feat(views): add JSP views, tag files, stylesheet and mini-cart}\\
Tests & \code{test: add automated checks and Verifications entry point}\\
Documentation & \code{docs: add README and project report}\\
\bottomrule
\end{tabularx}
\end{table}

